\documentclass[11pt]{article}
\usepackage[review]{acl}
\usepackage{times}
\usepackage{latexsym}
\usepackage{amsmath,amssymb,amsfonts}
\usepackage{graphicx}
\usepackage{booktabs}
\usepackage{multirow}
\usepackage{array}
\usepackage{microtype}
\usepackage{inconsolata}
\usepackage{hyperref}
\usepackage{placeins}
\providecommand{\citename}[1]{#1}

\title{A Benchmark for Sinhala-Script Language Identification: \\ Disambiguating Sinhala, Pali, and Sanskrit in Closed and Open World Contexts}

\author{Anonymous ACL submission}

\begin{document}
\maketitle

\begin{abstract}
State-of-the-art language identification (LangID) systems operate on an implicit orthographic assumption: that script uniquely identifies language. This assumption collapses catastrophically for the Sinhala script, which has historically been used to transcribe not only Sinhala, but also liturgical Pali and classical Sanskrit. Consequently, existing systems are structurally incapable of recognizing Sinhala-script Pali and Sanskrit, indiscriminately absorbing both into a single monolithic Sinhala label. In this paper, we present the first dedicated LangID benchmark and empirical study addressing this script-sharing challenge across two complementary research tiers. In \textbf{Phase 1 (Closed-World Benchmark)}, we evaluate seven baseline architectures trained strictly from scratch across four levels of input granularity (full sentences, 5-word, 3-word, and 1-word fragments). We discover that while full sentences saturate ($>0.994$ Macro-F1), single-word stress testing reveals a stark architectural divergence: probabilistic character $n$-gram models (Multinomial Naive Bayes: 0.8584 Macro-F1) and subword embeddings (fastText: 0.8027 Macro-F1) substantially outperform deep neural networks (Char-CNN: 0.7654; Char-BiGRU: 0.6916) and margin classifiers (Linear SVM: 0.6099) by exploiting high-likelihood inflectional affixes without requiring sentence-level syntax. In \textbf{Phase 2 (Open-World Benchmark)}, we scale the evaluation across three standardized hybrid benchmarks (FLORES+, CommonLID, WiLI-2018) spanning 11 languages. We demonstrate that naive fine-tuning induces catastrophic forgetting of global background languages, whereas targeted adaptation strategies—including Two-Stage Specialist Routing, LoRA rebalancing, and custom weight-preserving head extensions—eliminate script ambiguity while maintaining $0.92$--$0.97$ Macro-F1 globally. Finally, we uncover an inherent trade-off in massive multilingual LangID—the \textbf{Dialect Continuum Problem}—showing how fine-grained regional dialect heads in models like OpenLID-v2 induce fragmentation on classical languages.
\end{abstract}

\section{Introduction}
Language Identification (LangID) is an indispensable first stage in Natural Language Processing (NLP) pipelines, governing corpus curation, machine translation filtering, and information retrieval \cite{jauhiainen2019automatic,lui2012langid,joulin2017bag,kargaran2023glotlid}. While LangID is often considered an established problem for high-resource languages with distinctive orthographies, modern classifiers fundamentally struggle when multiple distinct languages share the exact same script \cite{caswell2020language,dunn2023commonlid,burchell2023openlid}. Most off-the-shelf classifiers rely heavily on character-block Unicode properties, falling victim to the \textit{Orthographic Fallacy}—the implicit inductive bias that a given script maps bijectively to a single dominant language.

This failure mode is particularly pronounced for the Sinhala script (Unicode block \texttt{U+0D80}--\texttt{U+0DFF}). In Sri Lanka and South Asian scholarly traditions, the Sinhala script has served for over two millennia as the primary written medium not only for \textbf{Sinhala} (an Indo-Aryan language spoken by over 17 million people), but also for \textbf{Pali} (the canonical liturgical language of Theravada Buddhism) and \textbf{Sanskrit} (the classical language of South Asian philosophy, medicine, and mathematics). The Buddhist monastic canon (\textit{Tipitaka}) was first committed to writing in the Sinhala script in the first century BCE at Aluvihara, Sri Lanka \cite{norman1983pali}, yielding an immense corpus of palm-leaf manuscripts and commentaries transcribed in Sinhala orthography.

Because existing planetary-scale LangID tools—such as fastText LID-176 \cite{joulin2017bag}, GlotLID v3 \cite{kargaran2023glotlid}, NLLB LID-218 \cite{nllb2022}, and OpenLID-v2 \cite{burchell2023openlid}—associate the Sinhala script exclusively with modern Sinhala (\texttt{sin} / \texttt{sin\_Sinh}), they exhibit near-zero recall on Sinhala-script Pali and Sanskrit. This systematic misclassification contaminates multilingual web crawls, corrupts machine translation pipelines, and impairs digital humanities scholarship on South Asian historical literature.

To resolve this challenge, this paper introduces a comprehensive benchmark and empirical investigation structured across two complementary phases:
\begin{enumerate}
    \item \textbf{Dedicated Multi-Language Target Corpus:} We construct a curated, group-stratified dataset of 74,318 sentences covering Sinhala, Pali, and Sanskrit written entirely in the Sinhala script, ensuring zero data leakage across document partitions.
    \item \textbf{Phase 1: Closed-World Baselines \& Morphological Stress Testing:} We train seven representative machine learning and neural architectures strictly from scratch. By subjecting them to consistent fragment stress testing (Full, 5-Word, 3-Word, and 1-Word inputs), we reveal the mechanisms of morphological discrimination under extreme contextual deprivation.
    \item \textbf{Phase 2: Open-World Benchmarking \& Catastrophic Forgetting Mitigation:} We construct three non-contaminated 11-language hybrid benchmarks based on FLORES+, CommonLID, and WiLI-2018. We benchmark six state-of-the-art foundation model families, demonstrating that naive fine-tuning causes severe catastrophic forgetting of pre-trained global languages.
    \item \textbf{Targeted Architectural Adaptation:} We engineer and validate five parameter-efficient and structural adaptation strategies (including Two-Stage Specialist Routing, LoRA with background rebalancing, and C++ weight-preserving head extension) that resolve script ambiguity while eliminating or minimizing degradation on global background languages (achieving 0.0\% degradation via Two-Stage Routing).
    
\end{enumerate}

\section{Dataset Construction and Taxonomy}

\subsection{Curated Target Corpus}
\label{sec:curated_target_corpus}
We compiled a standardized corpus from historical, canonical, and contemporary digital sources to benchmark language discrimination within the Sinhala script (see Appendix \ref{sec:appendix_data_provenance} for comprehensive provenance and download links):
\begin{itemize}
    \item \textbf{Sinhala (\texttt{Sinh-Sinh}, 32,263 sentences):} Curated from the SiPaKosa digital repository\footnote{\url{https://huggingface.co/datasets/RaniduG/SiPaKosa-Sent}} (Tipitaka-related literature, older books, and Buddhist biographical writing) and from the Sinhala side of a sentence-aligned Pali--Sinhala parallel corpus of canonical translations\footnote{\url{https://huggingface.co/datasets/sinhala-nlp/pali-sinhala}}, complemented by open digital libraries and contemporary journalistic publications.
    \item \textbf{Pali (\texttt{Pali-Sinh}, 29,317 sentences):} Extracted from the Pali side of the same sentence-aligned parallel corpus and from the SiDiaC-v.2.0 corpus\footnote{\url{https://github.com/NeviduJ/SiDiaC-v.2.0}} \cite{welgama2021sidiac}, which contributes verified digital editions of the Theravada \textit{Tipitaka}, its commentaries (\textit{Atthakatha}), sub-commentaries (\textit{Tika}), and Pali grammatical treatises \cite{norman1983pali}, all originally inscribed in Sinhala orthography.
    \item \textbf{Sanskrit (\texttt{San-Sinh}, 12,738 sentences):} Natively digitized Sanskrit in Sinhala script is scarce, so this class draws on three complementary sources. First, \textbf{SansinNT}\footnote{\url{https://ebible.org/details.php?id=sansin}} (7,953 sentences) supplies a historical Sanskrit New Testament translation natively rendered in the Sinhala script across 27 books. Second, classical \textit{kavya} and \textit{shastra} texts from the \textbf{Digital Corpus of Sanskrit} (DCS\footnote{\url{https://github.com/ambuda-org/dcs}}; 4,684 sentences across 12 works, e.g. \textit{Kumarasambhavam}, \textit{Meghadutam}, \textit{Saundaranandam}) \cite{hellwig2010dcs}, encoded in standard SLP1 Romanization, were algorithmically transliterated into the Sinhala script using the verified phonetic mappings of the \textit{Aksharamukha} Asian Script Converter \cite{rajan2023aksharamukha}. Third, a natively Sinhala-script Sanskrit devotional text (a \textit{kavaca} compendium of planetary and goddess invocations) from SiDiaC-v.2.0 \cite{welgama2021sidiac} contributes a further 101 sentences.
\end{itemize}

To prevent lexical memorization and optimistic evaluation bias, a strict \textbf{group-stratified split} was enforced. Sentences belonging to the same literary work, chapter, or commentary were assigned as atomic blocks into either training, validation, or testing partitions. The resulting unified target dataset comprises \textbf{74,318 sentences} (32,263 Sinhala / 43.4\%, 29,317 Pali / 39.4\%, and 12,738 Sanskrit / 17.1\%), partitioned as follows:
\begin{itemize}
    \item \textbf{Training Set (60,285 sentences / 81.1\%):} Used for feature extraction, model parameter optimization, and neural training (26,675 Sinhala, 23,490 Pali, 10,120 Sanskrit).
    \item \textbf{Validation Set (6,986 sentences / 9.4\%):} Reserved for loss monitoring, hyperparameter tuning, and early stopping (2,895 Sinhala, 2,800 Pali, 1,291 Sanskrit).
    \item \textbf{Held-Out Test Set (7,047 sentences / 9.5\%):} Strictly isolated from all training steps and used exclusively for final evaluation across full-sentence and fragmented benchmarks (2,693 Sinhala, 3,027 Pali, 1,327 Sanskrit).
\end{itemize}

\subsection{Phase 2 Hybrid Benchmarks and BCP-47 Tagging}
Evaluating LangID solely within a closed target world creates an artificial laboratory environment. In real-world deployment, an identification system must process target texts interspersed with dominant regional and global languages (e.g., English, Tamil, Hindi, Arabic).

Standard international benchmarks (FLORES+, CommonLID, WiLI-2018) universally assign text in the Sinhala script to a monolithic label (\texttt{sin\_Sinh}, \texttt{sin}, or \texttt{si}). To evaluate open-world discriminative capacity without test contamination, our taxonomy strictly adheres to the \textbf{BCP-47 language tag standard} \cite{phillips2009tags}, formatting each class as \texttt{[Language]-[Script]}:
\begin{itemize}
    \item \textbf{Target Classes (Sinhala Script):} \texttt{Sinh-Sinh}, \texttt{Pali-Sinh}, and \texttt{San-Sinh}.
    \item \textbf{Global Background Classes:} \texttt{San-Deva} (Devanagari Sanskrit), \texttt{Eng-Latn} (English), \texttt{Tam-Taml} (Tamil), \texttt{Hin-Deva} (Hindi), \texttt{Ben-Beng} (Bengali), \texttt{Ara-Arab} (Arabic), \texttt{Fre-Latn} (French), and \texttt{Ger-Latn} (German).
\end{itemize}

We constructed three hybrid evaluation benchmarks by isolating the official test partitions, replacing generic Sinhala rows with our 7,047 held-out target test sentences, and explicitly preserving Devanagari Sanskrit alongside the background languages:
\begin{enumerate}
    \item \textbf{FLORES+ Hybrid Benchmark (229,687 sentences):} Generic \texttt{sin\_Sinh} removed; contains 7,047 target test rows, 1,012 \texttt{san\_Deva} rows, and all FLORES+ background languages \cite{goyal2022flores}.
    \item \textbf{CommonLID Hybrid Benchmark (380,277 sentences):} Generic \texttt{sin} removed; contains 7,047 target test rows, 895 \texttt{san\_Deva} rows, and all CommonLID background languages \cite{dunn2023commonlid}.
    \item \textbf{WiLI-2018 Hybrid Benchmark (241,047 sentences):} Generic \texttt{si} removed; contains 7,047 target test rows, 1,000 \texttt{san\_Deva} rows, and all WiLI-2018 background languages \cite{thoma2018wili}.
\end{enumerate}
While these physical benchmark files retain the complete planetary-scale background (spanning over 200 languages and up to 380,277 sentences) to test open-world false-positive rejection, comparative benchmarking across all from-scratch baselines and foundation models is evaluated on the standardized \textbf{11-language core taxonomy} defined above (3 target Sinhala-script classes + 8 representative global background classes).

\subsection{Non-Contaminated Multi-Language Baseline Training Corpus}
\label{sec:multilang_train_corpus}
To train our seven representative from-scratch baseline architectures on the 11-language open-world task (Section \ref{sec:phase2_open_world}) without leaking evaluation data, we constructed a dedicated, uniformly balanced multi-language corpus of \textbf{110,000 sentences} (exactly 10,000 sentences per class across all 11 languages). Pre-trained foundation models already possess massive planetary-scale representations and were evaluated zero-shot or adapted via parameter-preserving strategies (Section \ref{sec:phase2_open_world}); in contrast, this dataset was engineered specifically to train the from-scratch baselines under clean, non-contaminated conditions. The dataset is partitioned into a 90\% training split ($N=99,000$; 9,000 sentences per class) and a 10\% validation split ($N=11,000$; 1,000 sentences per class), ensuring that every language holds an identical $9.09\%$ share to eliminate majority-class frequency bias.

To guarantee zero test-set contamination against the FLORES+, CommonLID, and WiLI-2018 benchmarks, data sources were curated under strict out-of-distribution isolation:
\begin{itemize}
    \item \textbf{Target Classes (\texttt{Sinh-Sinh}, \texttt{Pali-Sinh}, \texttt{San-Sinh}):} Sourced exclusively by sampling 10,000 sentences each from our dedicated target training split ($N=60,285$, Section \ref{sec:curated_target_corpus}), ensuring zero document or sentence overlap with the 7,047 held-out target test set.
    \item \textbf{Devanagari Sanskrit (\texttt{San-Deva}):} 10,000 classical Sanskrit sentences in Devanagari script were extracted from the open-source \textit{sanskrit\_classic} corpus on Hugging Face,\footnote{\url{https://huggingface.co/datasets/surajp/sanskrit_classic}} providing high-register epic and philosophical texts to force classifiers to differentiate language from orthography.
    \item \textbf{Tamil (\texttt{Tam-Taml}):} 10,000 contemporary Tamil sentences in native Tamil script were extracted from the official Wikimedia Tamil Wikipedia database dumps.\footnote{\url{https://dumps.wikimedia.org/tawiki/}}
    \item \textbf{Global Background Classes (\texttt{Eng-Latn}, \texttt{Hin-Deva}, \texttt{Ben-Beng}, \texttt{Ara-Arab}, \texttt{Fre-Latn}, \texttt{Ger-Latn}):} 10,000 clean sentences per language were curated from the open-source multilingual Tatoeba Project corpus.\footnote{\url{https://tatoeba.org/}}
\end{itemize}
All candidate sentences were filtered through a standardized quality-assurance pipeline enforcing a length constraint of 15--300 characters (3--60 words), collapsing redundant whitespace, and removing within-class duplicates prior to stratified sampling.

\section{Phase 1: Closed-World Baselines \& Fragment Stress Testing}

% \begin{table}[t]
% \centering
% \small
% \caption{Phase 1 Full Sentence Evaluation across Seven From-Scratch Baseline Models on the Held-Out Target Test Set ($N=7,047$). All models are trained exclusively on the target training split.}
% \label{tab:fullsent_baselines}
% \begin{tabular}{llcc}
% \toprule
% \textbf{Model} & \textbf{Paradigm} & \textbf{Accuracy} & \textbf{Macro-F1} \\
% \midrule
% Multinomial NB & Probabilistic ML & \textbf{0.9986} & \textbf{0.9985} \\
% Linear SVM & Margin Linear ML & 0.9982 & 0.9984 \\
% fastText-Scratch & Subword Embeddings & 0.9980 & 0.9980 \\
% Char $n$-gram + LogReg & Frequency NLP & 0.9977 & 0.9979 \\
% Char-BiGRU & Deep Recurrent & 0.9969 & 0.9970 \\
% Char-CNN & Deep 1D-CNN & 0.9960 & 0.9963 \\
% XGBoost & Gradient Boosting & 0.9946 & 0.9948 \\
% \bottomrule
% \end{tabular}
% \end{table}



\subsection{Baseline Architectural Setup}
We implemented seven representative architectures representing five distinct algorithmic paradigms, trained with zero prior external weights to establish an uncompromised closed-world baseline:
\begin{enumerate}
    \item \textbf{Multinomial Naive Bayes (MNB):} Character $n$-gram feature extraction (range $2$--$4$), TF-IDF weighted, with Laplace smoothing ($\alpha=1.0$) \cite{mccallum1998comparison}.
    \item \textbf{Linear Support Vector Machine (SVM):} Margin-based classifier trained with hinge loss and $L_2$ regularization on character $2$--$4$ gram TF-IDF vectors \cite{joachims1998text}.
    \item \textbf{fastText (Trained from Scratch):} Shallow two-layer architecture learning 256-dimensional subword character $n$-gram representations ($n \in [2, 5]$) trained with hierarchical softmax ($lr=0.5$, 15 epochs) \cite{joulin2017bag}.
    \item \textbf{Char $n$-gram + Logistic Regression:} Multinomial logistic regression over character $2$--$4$ grams with $L_2$ penalty and SAGA solver implemented in scikit-learn \cite{pedregosa2011scikit}.
    \item \textbf{Char-CNN:} Deep 1D convolutional neural network processing one-hot character index embeddings (vocab size 256, embedding dim 64) with three parallel convolution filters (kernel sizes 3, 5, 7; 128 channels each), global max pooling, dropout ($p=0.3$), and dense output \cite{zhang2015character}.
    \item \textbf{Char-BiGRU:} Two-layer bidirectional Gated Recurrent Unit (hidden dimension 128) over character embeddings with final hidden-state concatenation and linear projection \cite{cho2014learning}.
    \item \textbf{XGBoost:} Gradient boosted tree ensemble trained on sub-sampled character $n$-gram feature vectors (max depth 6, 200 estimators, learning rate 0.1) \cite{chen2016xgboost}.
\end{enumerate}

\begin{table*}[t]
\centering
\small
\caption{Phase 1 Consistent Fragment Stress Testing across all Seven From-Scratch Baseline Models. All models are trained strictly on the target training set ($N=60,285$) and evaluated on the held-out target test set ($N=7,047$) across full sentences, 5-word, 3-word, and 1-word inputs. Best results per fragment category are highlighted in bold.}
\label{tab:frag_stress_consistent}
\begin{tabular}{ll cccccccc}
\toprule
\multirow{2}{*}{\textbf{Model}} & \multirow{2}{*}{\textbf{Architecture Paradigm}} & \multicolumn{2}{c}{\textbf{Full Sentence}} & \multicolumn{2}{c}{\textbf{5-Word}} & \multicolumn{2}{c}{\textbf{3-Word}} & \multicolumn{2}{c}{\textbf{1-Word}} \\
\cmidrule(lr){3-4} \cmidrule(lr){5-6} \cmidrule(lr){7-8} \cmidrule(lr){9-10}
& & \textbf{Acc} & \textbf{M-F1} & \textbf{Acc} & \textbf{M-F1} & \textbf{Acc} & \textbf{M-F1} & \textbf{Acc} & \textbf{M-F1} \\
\midrule
Multinomial NB & Probabilistic ML & \textbf{0.9986} & \textbf{0.9985} & \textbf{0.9973} & \textbf{0.9972} & \textbf{0.9899} & \textbf{0.9885} & \textbf{0.8714} & \textbf{0.8584} \\
fastText-Scratch & Subword Embeddings & 0.9980 & 0.9980 & 0.9938 & 0.9934 & 0.9746 & 0.9702 & 0.8344 & 0.8027 \\
Linear SVM & Margin Linear ML & 0.9982 & 0.9984 & 0.9891 & 0.9880 & 0.9370 & 0.9309 & 0.6503 & 0.6099 \\
Char $n$-gram + LogReg & Frequency NLP & 0.9977 & 0.9979 & 0.9837 & 0.9800 & 0.9361 & 0.9168 & 0.6985 & 0.6223 \\
Char-CNN & Deep 1D-CNN & 0.9960 & 0.9963 & 0.9842 & 0.9821 & 0.9519 & 0.9434 & 0.7896 & 0.7654 \\
Char-BiGRU & Deep Recurrent & 0.9969 & 0.9970 & 0.9705 & 0.9630 & 0.9241 & 0.9015 & 0.7617 & 0.6916 \\
XGBoost & Gradient Boosting & 0.9946 & 0.9948 & 0.9679 & 0.9632 & 0.8886 & 0.8682 & 0.5509 & 0.4516 \\
\bottomrule
\end{tabular}
\end{table*}

\subsection{Full Sentence Saturation}
Table \ref{tab:frag_stress_consistent} reports full sentence evaluation on the held-out target test set. Across all seven architectures, performance reaches effective saturation, with Macro-F1 scores ranging between $0.9948$ (XGBoost) and $0.9985$ (Multinomial NB). When complete sentence syntax, case endings, and lexical co-occurrences are available, even shallow linear classifiers distinguish Sinhala, Pali, and Sanskrit with near-zero error, yielding virtually flawless diagonal predictions across probabilistic, linear, and neural baselines.

\subsection{Consistent Fragment Stress Testing}
Full sentences represent an idealized testing scenario. In historical palm-leaf manuscripts, epigraphy, digital humanities search queries, and noisy OCR pipelines, input text is frequently degraded into isolated phrases or individual words.

To measure true morphological discriminative capacity, we executed a rigorous fragment stress test. From each of the 7,047 held-out test sentences, contiguous fragments of 5 words, 3 words, and exactly 1 word were extracted. Table \ref{tab:frag_stress_consistent} provides the complete, consistent evaluation across all seven architectures.

As input length contracts, a pronounced architectural divergence emerges:
\begin{itemize}
    \item \textbf{Resilience of Probabilistic and Subword Architectures:} On single-word inputs, Multinomial NB retains an impressive \textbf{0.8714 Accuracy and 0.8584 Macro-F1}, while fastText achieves \textbf{0.8344 Accuracy and 0.8027 Macro-F1}.
    \item \textbf{Degradation of Deep Neural Networks:} Char-CNN drops to 0.7654 Macro-F1, while Char-BiGRU degrades to 0.6916 Macro-F1.
    \item \textbf{Collapse of Margin and Tree Classifiers:} Linear SVM drops precipitously to 0.6099 Macro-F1, Char+LogReg degrades to 0.6223 Macro-F1, and XGBoost experiences severe breakdown, achieving only 0.4516 Macro-F1.
\end{itemize}

\section{Morphological Analysis: Why Subwords Outperform Deep Networks}

The empirical supremacy of Multinomial Naive Bayes and subword fastText on 1-word fragments illuminates the linguistic structure of the target languages. Sinhala, Pali, and Sanskrit share extensive cognate roots due to their common Indo-Aryan ancestry, but diverge fundamentally in their inflectional morphology:
\begin{itemize}
    \item \textbf{Pali Morphological Signatures:} Pali preserves distinct Middle Indo-Aryan case markers that do not occur in classical Sanskrit or modern Sinhala. In particular, the genitive/dative singular suffix \textit{-ssa} (e.g., \textit{dhammassa}), ablative singular suffixes \textit{-smā} / \textit{-mhā}, and verbal third-person plural suffixes \textit{-enti} / \textit{-onti} serve as deterministic diagnostic character sequences.
    \item \textbf{Sanskrit Morphological Signatures:} Classical Sanskrit retains complex consonant conjuncts (\textit{saṃyuktākṣara}) absent in Pali's assimilated phonology (e.g., Sanskrit \textit{kṣetra} vs. Pali \textit{khetta}), along with explicit visarga markers (\textit{ḥ}, Sinhala glyph \texttt{U+0D83}) and virama-suppressed coda stops.
    \item \textbf{Sinhala Morphological Signatures:} Modern Sinhala exhibits distinctive vowel signs, notably the low front vowels \textit{æ} (\texttt{U+0DD0}) and \textit{ǣ} (\texttt{U+0DD1}); also radical case-ending simplification.
\end{itemize}

Probabilistic Naive Bayes operates as a cumulative likelihood accumulator over character $n$-grams. When presented with an isolated word containing the subword string \textit{-ssa}, the conditional log-likelihood $\log P(\text{\textit{-ssa}} \mid \text{Pali})$ decisively outvotes alternative classes, irrespective of sentence context. fastText operates via a similar mechanism, summing learned subword vector representations directly into a document vector \cite{mikolov2013distributed}.

Conversely, deep sequence models (CNNs and BiGRUs) optimize filters for composite sequence structures and token transitions learned across full sentence lengths. When stripped of sentence context and forced onto short, zero-padded fragments, the spatial feature activations are disrupted. Linear SVMs and XGBoost, in turn, rely on sparse hyperplanes whose margin boundaries are heavily calibrated on sentence-level feature density, leading to severe misclassifications on isolated root words.

\FloatBarrier
\section{Phase 2: Open-World SOTA Benchmarking \& Catastrophic Forgetting}
\label{sec:phase2_open_world}

While Phase 1 proves that from-scratch models can resolve target script ambiguity in a closed world, deploying them in open-world production pipelines introduces the catastrophic forgetting dilemma. If a model is trained exclusively on target languages, it loses the ability to recognize global languages. Conversely, if off-the-shelf foundation models are deployed zero-shot, they fail on the target languages.

\subsection{Zero-Shot Foundation Model Collapse}
We evaluated six state-of-the-art foundation model families in a zero-shot regime on our hybrid benchmarks: multilingual transformers (XLM-RoBERTa Base \cite{conneau2020unsupervised}, following the multilingual masked language modeling paradigm of mBERT \cite{devlin2019bert}), ConLID \cite{dunn2023commonlid}, fastText LID-176 \cite{joulin2017bag}, GlotLID v3 \cite{kargaran2023glotlid}, NLLB LID-218 \cite{nllb2022}, and OpenLID-v2 \cite{burchell2023openlid}.

As detailed in Appendix \ref{sec:appendix_tables} (Tables \ref{tab:flores_full}, \ref{tab:common_full}, and \ref{tab:wili_full}), every stock model achieves \textbf{0.00\% F1 on Pali-Sinh and San-Sinh}. Because stock classification heads provide only a single generic Sinhala label (\texttt{sin} or \texttt{sin\_Sinh}), all Sinhala-script inputs are funneled into that class, depressing target recall to zero and yielding an aggregate target Macro-F1 of only $\approx 0.18$.

\subsection{Targeted Adaptation Strategies}
Direct, unconstrained fine-tuning on a narrow target dataset causes catastrophic forgetting, completely eroding pre-trained representations for global languages \cite{mccloskey1989catastrophic,french1999catastrophic,kirkpatrick2017overcoming}. To achieve high target precision while guaranteeing background retention, we engineered five architectural adaptation strategies:
\begin{enumerate}
    \item \textbf{Two-Stage Specialist Routing (OpenLID-v2 \& fastText LID-176):} A frozen pre-trained Stage 1 router screens incoming text. Non-Sinhala scripts are classified directly by the global model. When the Sinhala script is detected (\texttt{sin\_Sinh} / \texttt{sin}), text is diverted to an isolated Stage 2 specialist model fine-tuned for fine-grained 3-way discrimination (\texttt{Sinh-Sinh}, \texttt{Pali-Sinh}, \texttt{San-Sinh}). By design, this hierarchical routing guarantees \textbf{0.0\% performance degradation on all non-Sinhala languages}.
    \item \textbf{LoRA Parameter-Efficient Adaptation with Rebalancing (XLM-R Base):} Low-Rank Adaptation (LoRA; rank $r=8$, $\alpha=16$) matrices were injected into the query, key, and value attention projections of XLM-R Base \cite{hu2022lora}. During fine-tuning, training batches were dynamically rebalanced with background language exemplars from our non-contaminated corpus, anchoring the multilingual shared representations.
    \item \textbf{Weight-Preserving Head Extension \& C++ Continuation Training (NLLB LID-218):} A custom C++ compilation patch (\texttt{extend-labels}) was developed for fastText to expand the model's output projection matrix from 218 to 220 labels. Existing weight rows were preserved verbatim, while newly added label vectors were initialized and trained via continuation training (\texttt{finetune-loaded}) with linear learning rate decay ($lr=0.001$, 3 epochs).
    \item \textbf{In-Place Head Extension with Xavier Initialization (ConLID):} The final linear projection head of ConLID was expanded, preserving original pre-trained weights while initializing the new target heads with Xavier uniform initialization. Optimization utilized AdamW ($lr=1e-4$) and validation-loss early stopping.
    \item \textbf{Pretrained Vector Transfer (GlotLID v3):} Subword representations from GlotLID's pre-trained 2,100-language embedding space were transferred to initialize model adaptation, paired with conservative learning schedules ($lr=0.005$) to prevent weight drift.
\end{enumerate}

\subsection{Benchmark Performance Overview}

\begin{table}[!htbp]
\centering
\footnotesize
\setlength{\tabcolsep}{2.5pt}
\caption{Overall Macro-F1 across three Phase 2 Hybrid Benchmarks for Trained Baselines and Adapted Foundation Models across 11 languages (3 target + 8 global background).}
\label{tab:sota_overall_summary}
\resizebox{\columnwidth}{!}{%
\begin{tabular}{lccc}
\toprule
\textbf{Model Family} & \textbf{FLORES+} & \textbf{CommonLID} & \textbf{WiLI-18} \\
\midrule
\multicolumn{4}{l}{\textit{\textbf{From-Scratch Baselines} (Trained on Uniform 110k)}} \\
Multinomial NB & 0.9332 & 0.9342 & \textbf{0.9785} \\
Linear SVM & 0.9342 & 0.8895 & 0.9679 \\
Char $n$-gram + LogReg & 0.9298 & 0.8907 & 0.9750 \\
fastText (from scratch) & 0.9281 & 0.8758 & 0.9584 \\
XGBoost & 0.9181 & 0.8734 & 0.9215 \\
Char-CNN (1D-CNN) & 0.9177 & 0.8398 & 0.8790 \\
Char-BiGRU (Recurrent) & 0.9111 & 0.8166 & 0.8770 \\
\midrule
\multicolumn{4}{l}{\textit{\textbf{Adapted Pretrained Foundation Models}}} \\
NLLB LID-218 (C++ Patch) & \textbf{0.9530} & 0.9501 & 0.9675 \\
GlotLID v3 (Vec Transfer) & 0.9465 & 0.9527 & 0.9570 \\
ConLID (Head Ext.) & 0.9329 & 0.9261 & 0.9688 \\
fastText LID-176 (2-Stage) & 0.9276 & 0.9645 & 0.9745 \\
XLM-R Base (LoRA) & 0.9267 & \textbf{0.9737} & 0.8880 \\
OpenLID-v2 (2-Stage) & 0.7759 & 0.7913 & 0.7782 \\
\bottomrule
\end{tabular}%
}
\end{table}

Table \ref{tab:sota_overall_summary} presents the comprehensive Macro-F1 performance across all three hybrid benchmarks for both our seven from-scratch baseline models (trained on the uniform 110,000-sentence dataset) and the six adapted foundation models. 

A clear empirical dichotomy emerges between from-scratch baselines and adapted foundation models:
\begin{itemize}
    \item \textbf{Global Generalization Advantage of Adapted Foundation Models:} When evaluated on diverse, real-world web text, adapted foundation models achieve the highest peak Macro-F1 scores, notably NLLB LID-218 on FLORES+ (\textbf{0.9530}) and XLM-R Base with LoRA (\textbf{0.9737}) and fastText LID-176 2-Stage (\textbf{0.9645}) on CommonLID. This performance advantage stems from their multi-billion-token pre-training across hundreds of languages, which supplies rich, resilient lexicons for global background classes (e.g., English, French, Arabic, and Tamil) that overcome noisy orthographic variance on the open web where 10,000-sentence scratch baselines show slight degradation.
    \item \textbf{Computational Efficiency of From-Scratch Baselines:} Conversely, lightweight probabilistic and margin baselines exhibit remarkable competitive fidelity without requiring massive pretraining: Multinomial Naive Bayes attains \textbf{0.9785 Macro-F1 on WiLI-2018} (outperforming all foundation models) and \textbf{0.9342 on CommonLID}, while Linear SVM attains \textbf{0.9342 on FLORES+}. The un-pretrained deep learning baselines (Char-CNN at 0.8398 and Char-BiGRU at 0.8166 on CommonLID) demonstrate greater sensitivity to noisy web distributions compared to frequency-based subwords.
\end{itemize}

Crucially, while stock foundation models collapse zero-shot ($0.0\%$ target recall), targeted foundation adaptation completely eliminates the orthographic failure mode, lifting target-class F1 to 0.95--0.99 while maintaining pristine background retention across all 11 evaluation classes.

\FloatBarrier
\section{Discussion: The Dialect Continuum Problem and Devanagari Fragmentation}

A critical empirical discovery in Table \ref{tab:sota_overall_summary} is the stark divergence between OpenLID-v2 (Macro-F1 $0.7759$--$0.7913$) and other models like fastText LID-176 ($0.9276$--$0.9745$) and NLLB LID-218 ($0.9501$--$0.9675$).

Detailed inspection of the confusion matrices reveals that this performance gap stems entirely from the treatment of Devanagari Sanskrit (\texttt{San-Deva}). In FLORES+, OpenLID-v2 achieves a dismal \textbf{0.0350 F1 on \texttt{San-Deva}}, whereas fastText LID-176 achieves \textbf{0.9594 F1}.

This discrepancy highlights a fundamental manifestation of the \textbf{Dialect Continuum Problem} in massive multilingual LangID. Across continuous linguistic zones—such as the Indo-Aryan dialect continuum connecting Sanskrit, Prakrits, and modern regional varieties—languages share dense lexical, phonological, and subword inventories without sharp geographical boundaries. OpenLID-v2 was trained with an expansive label space designed to identify granular regional dialects across South Asia. When presented with classical Devanagari Sanskrit text, OpenLID-v2 fragments the probability mass across multiple genetically descended dialect heads, misclassifying 68.1\% of Sanskrit sentences as Bhojpuri (\texttt{bho\_Deva}), 15.1\% as Maithili (\texttt{mai\_Deva}), and 43.3\% of Hindi sentences as Kashmiri (\texttt{kas\_Deva}). In contrast, fastText LID-176 and NLLB LID-218 employ a unified classical language tag (\texttt{sa} / \texttt{san\_Deva}), avoiding continuous boundary leakage and preserving high classification fidelity. This finding demonstrates that for classical literary languages, ultra-granular dialect heads can actively undermine identification accuracy by exacerbating dialect continuum ambiguity.

\section{Practical Recommendations}
Based on our empirical findings across closed and open benchmarks, we provide the following practical recommendations for NLP practitioners:
\begin{enumerate}
    \item \textbf{Short-Text \& OCR Applications:} For historical manuscript OCR and search queries consisting of isolated words or phrases, practitioners should deploy \textbf{Multinomial Naive Bayes or shallow fastText classifiers} over character $n$-grams. Massive neural transformers are ill-suited for short-text LangID without syntactic context.
    \item \textbf{Web-Scale Multilingual Pipelines:} For planetary-scale web crawling (e.g., Common Crawl), practitioners should adopt \textbf{Two-Stage Specialist Routing}. This modular approach eliminates script ambiguity with mathematical certainty of zero degradation on pre-trained background languages.
    \item \textbf{Fine-Tuning Foundation Transformers:} When adapting models like XLM-R, practitioners must employ \textbf{parameter-efficient LoRA adapters combined with background replay rebalancing} to avoid catastrophic forgetting.
    \item \textbf{Accuracy vs. Compute Trade-Off (Adapted SOTA vs. Baselines):} When GPU compute is available and maximum global generalization across noisy, unpredictable web text is required, \textbf{adapted foundation models} (e.g., XLM-R Base with LoRA or fastText LID-176 2-Stage) are the recommended choice due to their superior pre-trained background robustness. Conversely, for CPU-bound or edge deployment, \textbf{from-scratch character $n$-gram baselines} deliver within 1.5--4\% of SOTA Macro-F1 with zero GPU footprint and sub-millisecond inference.
\end{enumerate}

\section{Limitations}
First, due to the scarcity of natively digitized Sanskrit in Sinhala script, roughly 37\% of our Sanskrit target corpus (4,684 of 12,738 sentences, drawn from the Digital Corpus of Sanskrit) relies on algorithmic transliteration from SLP1 via Aksharamukha; the remainder (SansinNT and SiDiaC-v.2.0) is natively Sinhala-script. While phonetically deterministic, modern digital transliteration does not fully reflect historical scribal ligatures. Second, heavily degraded palm-leaf manuscripts containing archaic Sinhala glyphs and non-standard orthography require dedicated optical preprocessing before language identification can be performed.

\section{Conclusion}
This paper presents the first systematic benchmark and empirical investigation of Sinhala-Script Language Identification across Sinhala, Pali, and Sanskrit. In Phase 1, we demonstrate that probabilistic subword models achieve superior resilience on short text fragments by capturing diagnostic inflectional affixes. In Phase 2, we demonstrate that targeted adaptation strategies resolve the catastrophic forgetting dilemma: adapted foundation models leverage their massive pre-trained representations to achieve peak global robustness on noisy web data (up to 0.9737 Macro-F1), while lightweight from-scratch baselines offer competitive efficiency (up to 0.9785 Macro-F1) on clean domains. Our datasets, code, and adapted model weights are released to advance low-resource and digital humanities NLP.

% Entries for the entire Anthology, followed by custom entries
\bibliography{custom}

\appendix
\section{Detailed Hybrid Benchmark Results Across 11 Languages}
\label{sec:appendix_tables}

Tables \ref{tab:flores_full}, \ref{tab:common_full}, and \ref{tab:wili_full} provide the complete per-class Macro-F1 evaluation on FLORES+, CommonLID, and WiLI-2018, comparing zero-shot foundation models against adapted fine-tuned models across all 11 target and global background languages.

\begin{table*}[p]
\centering
\scriptsize
\setlength{\tabcolsep}{2.2pt}
\renewcommand{\arraystretch}{0.90}
\caption{FLORES+ Hybrid Benchmark Performance: Zero-Shot Baseline vs. Fine-Tuned Models across all 11 target and global background languages.}
\label{tab:flores_full}
\begin{tabular}{l ccc cccccccc c}
\toprule
\multicolumn{13}{c}{\textbf{Zero-Shot Baseline Results}} \\
\midrule
\textbf{Model} & \textbf{Sinh-Sinh} & \textbf{Pali-Sinh} & \textbf{San-Sinh} & \textbf{San-Deva} & \textbf{Eng-Latn} & \textbf{Tam-Taml} & \textbf{Hin-Deva} & \textbf{Ben-Beng} & \textbf{Ara-Arab} & \textbf{Fre-Latn} & \textbf{Ger-Latn} & \textbf{Macro F1} \\
\midrule
XLM-R Base & -- & -- & -- & -- & 0.9985 & -- & 0.1835 & -- & 0.6671 & 0.9990 & 1.0000 & -- \\
ConLID & 0.7916 & -- & -- & 0.9985 & 0.9960 & 0.9995 & 0.9970 & 0.9995 & 0.3733 & 1.0000 & 0.9970 & -- \\
fastText LID-176 & 0.7912 & -- & -- & 0.9594 & 0.7213 & 1.0000 & 0.9629 & 1.0000 & 0.6667 & 0.9873 & 0.9897 & -- \\
GlotLID v3 & 0.7914 & -- & -- & 0.9950 & 1.0000 & 1.0000 & 0.9926 & 0.9995 & 0.6054 & 1.0000 & 1.0000 & -- \\
NLLB LID-218 & 0.7912 & -- & -- & 0.9895 & 0.9840 & 1.0000 & 0.9912 & 1.0000 & 0.6664 & 0.9995 & 1.0000 & -- \\
OpenLID-v2 & 0.7912 & -- & -- & 0.0350 & 0.9995 & 1.0000 & 0.0361 & 0.9970 & 0.5579 & 0.9975 & 1.0000 & -- \\
\midrule
\multicolumn{13}{c}{\textbf{Fine-Tuned Model Results}} \\
\midrule
\textbf{Model} & \textbf{Sinh-Sinh} & \textbf{Pali-Sinh} & \textbf{San-Sinh} & \textbf{San-Deva} & \textbf{Eng-Latn} & \textbf{Tam-Taml} & \textbf{Hin-Deva} & \textbf{Ben-Beng} & \textbf{Ara-Arab} & \textbf{Fre-Latn} & \textbf{Ger-Latn} & \textbf{Macro F1} \\
\midrule
XLM-R Base & 0.9967 & 0.9844 & 0.9974 & -- & 0.7734 & 0.9912 & -- & 0.9970 & 0.6667 & 0.8604 & 1.0000 & 0.9267 \\
ConLID & 0.9634 & 0.9705 & 0.9414 & 0.9985 & 0.9960 & 0.9995 & 0.9970 & 0.9995 & 0.3991 & 1.0000 & 0.9970 & 0.9329 \\
fastText LID-176 & 0.9611 & 0.9751 & 0.9805 & 0.9594 & 0.7213 & 1.0000 & 0.9629 & 1.0000 & 0.6667 & 0.9873 & 0.9892 & 0.9276 \\
GlotLID v3 & 0.9688 & 0.9566 & 0.9019 & 0.9955 & 1.0000 & 1.0000 & 0.9926 & 0.9995 & 0.5962 & 1.0000 & 1.0000 & 0.9465 \\
NLLB LID-218 & 0.9760 & 0.9526 & 0.9053 & 0.9915 & 1.0000 & 1.0000 & 0.9921 & 1.0000 & 0.6667 & 0.9995 & 0.9995 & 0.9530 \\
OpenLID-v2 & 0.9670 & 0.9743 & 0.9704 & 0.0350 & 0.9995 & 1.0000 & 0.0361 & 0.9970 & 0.5579 & 0.9975 & 1.0000 & 0.7759 \\
\bottomrule
\end{tabular}

\vspace{12pt}

\caption{CommonLID Hybrid Benchmark Performance: Zero-Shot Baseline vs. Fine-Tuned Models across all 11 target and global background languages.}
\label{tab:common_full}
\begin{tabular}{l ccc cccccccc c}
\toprule
\multicolumn{13}{c}{\textbf{Zero-Shot Baseline Results}} \\
\midrule
\textbf{Model} & \textbf{Sinh-Sinh} & \textbf{Pali-Sinh} & \textbf{San-Sinh} & \textbf{San-Deva} & \textbf{Eng-Latn} & \textbf{Tam-Taml} & \textbf{Hin-Deva} & \textbf{Ben-Beng} & \textbf{Ara-Arab} & \textbf{Fre-Latn} & \textbf{Ger-Latn} & \textbf{Macro F1} \\
\midrule
XLM-R Base & -- & -- & -- & -- & 0.9473 & -- & 0.4403 & -- & 0.9955 & 0.9483 & 0.9677 & -- \\
ConLID & 0.7916 & -- & -- & 0.9591 & 0.8500 & 0.9818 & 0.9443 & 0.9615 & 0.7664 & 0.8974 & 0.8953 & -- \\
fastText LID-176 & 0.7910 & -- & -- & 0.8886 & 0.9724 & 0.9643 & 0.9702 & 0.9936 & 0.9947 & 0.9447 & 0.9673 & -- \\
GlotLID v3 & 0.7914 & -- & -- & 0.9634 & 0.9255 & 0.9878 & 0.9713 & 0.9792 & 0.9631 & 0.9255 & 0.9215 & -- \\
NLLB LID-218 & 0.7912 & -- & -- & 0.9282 & 0.9240 & 0.9747 & 0.9640 & 0.9863 & 0.9923 & 0.9289 & 0.9410 & -- \\
OpenLID-v2 & 0.7912 & -- & -- & 0.1169 & 0.8973 & 0.9877 & 0.1749 & 0.9722 & 0.9473 & 0.9052 & 0.9239 & -- \\
\midrule
\multicolumn{13}{c}{\textbf{Fine-Tuned Model Results}} \\
\midrule
\textbf{Model} & \textbf{Sinh-Sinh} & \textbf{Pali-Sinh} & \textbf{San-Sinh} & \textbf{San-Deva} & \textbf{Eng-Latn} & \textbf{Tam-Taml} & \textbf{Hin-Deva} & \textbf{Ben-Beng} & \textbf{Ara-Arab} & \textbf{Fre-Latn} & \textbf{Ger-Latn} & \textbf{Macro F1} \\
\midrule
XLM-R Base & 0.9966 & 0.9901 & 0.9966 & -- & 0.9786 & 0.9310 & 0.9925 & 0.9865 & 0.9971 & 0.9521 & 0.9665 & 0.9737 \\
ConLID & 0.9634 & 0.9705 & 0.9414 & 0.9607 & 0.8488 & 0.9818 & 0.9415 & 0.9592 & 0.7966 & 0.8975 & 0.8888 & 0.9261 \\
fastText LID-176 & 0.9609 & 0.9751 & 0.9805 & 0.8886 & 0.9724 & 0.9643 & 0.9702 & 0.9936 & 0.9947 & 0.9447 & 0.9672 & 0.9645 \\
GlotLID v3 & 0.9689 & 0.9566 & 0.9019 & 0.9612 & 0.9270 & 0.9818 & 0.9646 & 0.9792 & 0.9603 & 0.9256 & 0.9211 & 0.9527 \\
NLLB LID-218 & 0.9625 & 0.9500 & 0.9053 & 0.9310 & 0.9072 & 0.9875 & 0.9588 & 0.9874 & 0.9919 & 0.9192 & 0.9407 & 0.9501 \\
OpenLID-v2 & 0.9670 & 0.9743 & 0.9704 & 0.1169 & 0.8973 & 0.9877 & 0.1749 & 0.9722 & 0.9473 & 0.9052 & 0.9239 & 0.7913 \\
\bottomrule
\end{tabular}

\vspace{12pt}

\caption{WiLI-2018 Hybrid Benchmark Performance: Zero-Shot Baseline vs. Fine-Tuned Models across all 11 target and global background languages.}
\label{tab:wili_full}
\begin{tabular}{l ccc cccccccc c}
\toprule
\multicolumn{13}{c}{\textbf{Zero-Shot Baseline Results}} \\
\midrule
\textbf{Model} & \textbf{Sinh-Sinh} & \textbf{Pali-Sinh} & \textbf{San-Sinh} & \textbf{San-Deva} & \textbf{Eng-Latn} & \textbf{Tam-Taml} & \textbf{Hin-Deva} & \textbf{Ben-Beng} & \textbf{Ara-Arab} & \textbf{Fre-Latn} & \textbf{Ger-Latn} & \textbf{Macro F1} \\
\midrule
XLM-R Base & -- & -- & -- & -- & 0.9527 & -- & 0.1811 & -- & -- & 0.9900 & 0.9899 & -- \\
ConLID & 0.7916 & -- & -- & 0.9924 & 0.9439 & 0.9945 & 0.9899 & 0.9435 & -- & 0.9819 & 0.9707 & -- \\
fastText LID-176 & 0.7912 & -- & -- & 0.9904 & 0.9385 & 0.9950 & 0.9889 & 0.9529 & -- & 0.9774 & 0.9854 & -- \\
GlotLID v3 & 0.7914 & -- & -- & 0.9914 & 0.9377 & 0.9950 & 0.9173 & 0.9451 & -- & 0.9880 & 0.9754 & -- \\
NLLB LID-218 & 0.7912 & -- & -- & 0.9909 & 0.9301 & 0.9940 & 0.9889 & 0.9429 & -- & 0.9900 & 0.9904 & -- \\
OpenLID-v2 & 0.7912 & -- & -- & 0.0276 & 0.9427 & 0.9940 & 0.0134 & 0.9424 & -- & 0.9687 & 0.9817 & -- \\
\midrule
\multicolumn{13}{c}{\textbf{Fine-Tuned Model Results}} \\
\midrule
\textbf{Model} & \textbf{Sinh-Sinh} & \textbf{Pali-Sinh} & \textbf{San-Sinh} & \textbf{San-Deva} & \textbf{Eng-Latn} & \textbf{Tam-Taml} & \textbf{Hin-Deva} & \textbf{Ben-Beng} & \textbf{Ara-Arab} & \textbf{Fre-Latn} & \textbf{Ger-Latn} & \textbf{Macro F1} \\
\midrule
XLM-R Base & 0.9967 & 0.9957 & 0.9970 & -- & 0.9569 & 0.9950 & 0.9909 & 0.9701 & -- & 0.9886 & 0.9894 & 0.8880 \\
ConLID & 0.9634 & 0.9705 & 0.9414 & 0.9924 & 0.9433 & 0.9945 & 0.9899 & 0.9435 & -- & 0.9809 & 0.9680 & 0.9688 \\
fastText LID-176 & 0.9611 & 0.9751 & 0.9805 & 0.9904 & 0.9385 & 0.9950 & 0.9889 & 0.9529 & -- & 0.9774 & 0.9850 & 0.9745 \\
GlotLID v3 & 0.9688 & 0.9566 & 0.9019 & 0.9914 & 0.9367 & 0.9950 & 0.9126 & 0.9451 & -- & 0.9875 & 0.9744 & 0.9570 \\
NLLB LID-218 & 0.9760 & 0.9526 & 0.9053 & 0.9909 & 0.9455 & 0.9940 & 0.9879 & 0.9429 & -- & 0.9900 & 0.9899 & 0.9675 \\
OpenLID-v2 & 0.9670 & 0.9743 & 0.9704 & 0.0276 & 0.9427 & 0.9940 & 0.0134 & 0.9424 & -- & 0.9687 & 0.9817 & 0.7782 \\
\bottomrule
\end{tabular}
\end{table*}

\section{Dataset Provenance and Source Corpus Breakdown}
\label{sec:appendix_data_provenance}

To ensure complete scientific transparency and reproducibility, Table \ref{tab:dataset_provenance} summarizes the provenance, orthographic representation, textual domains, sentence distributions, and public repository links for all primary sources assembled into our 74,318-sentence target benchmark.

\begin{table*}[t]
\centering
\small
\setlength{\tabcolsep}{5pt}
\caption{Target Corpus Provenance and Source Breakdown for Sinhala, Pali, and Sanskrit in Sinhala Script ($N=74,318$). All source repositories are publicly accessible.}
\label{tab:dataset_provenance}
\resizebox{\textwidth}{!}{%
\begin{tabular}{llrlll}
\toprule
\textbf{Language Class} & \textbf{Primary Source / Repository} & \textbf{Sentences} & \textbf{Orthography} & \textbf{Genre / Domain} & \textbf{Public Access Repository} \\
\midrule
\multirow{2}{*}{Sinhala (\texttt{Sinh-Sinh})} 
& SiPaKosa-Sent & 20,418 & Native Sinhala & Buddhist Commentary \& History & \href{https://huggingface.co/datasets/RaniduG/SiPaKosa-Sent}{\texttt{hf.co/datasets/RaniduG/SiPaKosa-Sent}} \\
& Pali--Sinhala Parallel & 11,845 & Native Sinhala & Canonical Translation Exegesis & \href{https://huggingface.co/datasets/sinhala-nlp/pali-sinhala}{\texttt{hf.co/datasets/sinhala-nlp/pali-sinhala}} \\
\midrule
\multirow{2}{*}{Pali (\texttt{Pali-Sinh})} 
& SiDiaC-v.2.0 & 17,472 & Native Sinhala & Canonical Tipitaka \& Atthakatha & \href{https://github.com/NeviduJ/SiDiaC-v.2.0}{\texttt{github.com/NeviduJ/SiDiaC-v.2.0}} \\
& Pali--Sinhala Parallel & 11,845 & Native Sinhala & Canonical Sutta Verses & \href{https://huggingface.co/datasets/sinhala-nlp/pali-sinhala}{\texttt{hf.co/datasets/sinhala-nlp/pali-sinhala}} \\
\midrule
\multirow{3}{*}{Sanskrit (\texttt{San-Sinh})} 
& SansinNT & 7,953 & Native Sinhala & Sanskrit New Testament (27 Books) & \href{https://ebible.org/details.php?id=sansin}{\texttt{ebible.org/details.php?id=sansin}} \\
& DCS (Kalidasa, Asvaghosa) & 4,684 & Transliterated & Classical Kavya \& Shastra & \href{https://github.com/ambuda-org/dcs}{\texttt{github.com/ambuda-org/dcs}} \\
& SiDiaC (Kavacha Sangrahaya) & 101 & Native Sinhala & Protective Mantras \& Stotras & \href{https://github.com/NeviduJ/SiDiaC-v.2.0}{\texttt{github.com/NeviduJ/SiDiaC-v.2.0}} \\
\midrule
\textbf{Total Target Corpus} & & \textbf{74,318} & & & \\
\bottomrule
\end{tabular}%
}
\end{table*}

\end{document}
